\begin{table}[!t]
    \centering \footnotesize
    \begin{tabular}{p{.95\columnwidth}}
        \texttt{New egg cells form via a specialized kind of cell division called called meiosis, and will pause at key stages in this process before continuing their development. One of these pauses occurs before the egg cell is fertilized. At fertilization, the egg cell becomes “activated”, development resumes, and it starts forming into an embryo. 
        Molecules deposited in the egg cell when it originally formed are used to control these earliest stages of embryonic development. These molecules include messenger RNA molecules (mRNAs for short), which can be “translated” to build proteins. In fruit flies, an enzyme called PNG kinase regulates the translation of hundreds of mRNA molecules during the period after the pause, when the maturing egg cell is activated and the embryo begins to develop. 
        [...]
        % It is not well understood what activates and inactivates the kinase to limit its activity to this period of time. However, it was known that a protein called GNU was needed to bind to the PNG kinase to make it active. CyclinB/CDK1 is another kinase, and in contrast to PNG it is highly active when the egg cell is paused. When the egg cell is activated for embryonic development, the levels of this second kinase drop sharply and meiosis is completed. Like all kinases, CyclinB/CDK1 attaches phosphate groups onto other molecules, and Hara et al. now show that CyclinB/CDK1 can modify the GNU protein in this way. The added phosphate groups prevent GNU from binding to the PNG kinase, meaning that the high levels of CyclinB/CDK1 during the pause stop GNU from activating the PNG kinase. However, when the egg cell is activated, the level of CyclinB/CDK1 declines so that there are not enough of these molecules to add phosphates onto GNU. This leaves GNU free to activate the PNG kinase, allowing this kinase to control the translation of mRNA molecules. Furthermore, the activity of PNG kinase leads to the destruction of GNU, and this feedback loop limits this kinase’s activity to the narrow window of time in which it is needed. 
        The fruit fly is the second example of an animal in which the activity of a kinase essential for embryonic development has been linked to the completion of meiosis (the other being the roundworm Caenorhabditis elegans). The use of this strategy in two such different animals suggests that it may also be common to many other animals, including humans. Further investigation is now needed to determine if this is indeed the case.}
    \end{tabular}
    \caption{Excerpt of a summary from eLife. Coleman score = 10.5 (10-11th grade proficiency), LM score = 2}
    \label{tab:disagreement-example2}
\end{table}